\documentclass[11pt]{article}
\usepackage[a4paper,margin=1in]{geometry}
\usepackage{amsmath,amssymb,amsthm}
\usepackage{booktabs}
\usepackage{graphicx}
\usepackage{xcolor}
\usepackage{hyperref}
\hypersetup{colorlinks=true,linkcolor=blue,urlcolor=blue,citecolor=blue}
\newcommand{\rt}[1]{\texttt{#1}}
\newcommand{\fn}[1]{\texttt{#1}}
\newcommand{\wcc}{\mathrm{wcc}}
\newcommand{\scc}{\mathrm{scc}}
\newcommand{\ifgraphic}[2]{\IfFileExists{#1}{\includegraphics[width=#2]{#1}}%
  {\fbox{\parbox{#2}{\centering\small figure \texttt{\detokenize{#1}} pending}}}}

\title{\textbf{Report R13 --- nested fragmentation: can a rule fragment
twice?}\\
\large Krylov sector analysis of quantum cellular automata}
\author{Tier 1 analysis}
\date{2026-08-04}

\begin{document}
\maketitle

\begin{abstract}
A rule fragments once when its number of sectors grows exponentially in $N$. It
fragments a \emph{second} time when a single sector contains exponentially many
attractors. Only thirteen rules can do the second thing at all --- every other
rule of the $256$ has exactly one attractor per sector (R9~\S7.3) --- so the
question is decidable by inspecting thirteen sequences. The answer is
\textbf{partly}. Eight rules fragment exponentially \emph{inside} a sector while
their sector count stays constant or linear, and that half is settled by exact
integer recurrences rather than fits. Three constants appear, all
``Narayana-like'': the supergolden $\psi=1.465571$, the plastic number
$\rho=1.324718$, and the root $1.380278$ of $x^4=x^3+1$. The distribution of
attractors across sectors for rules $44$ and $100$ is, exactly, a run of
\emph{consecutive Narayana numbers}. The fully nested case --- exponential both
ways --- has two candidates, rules $29$ and $71$, and \textbf{both are left
open}. For rule $29$, on the even-$N$ branch up to $N=20$ a quadratic and an
exponential of base $1.12$ fit its in-sector count to within $6\%$ with
essentially equal BIC, and they do not separate by a full attractor until
$N\simeq26$, which is $6.7\times10^{7}$ states. Its near-twin rule $71$ is exactly linear inside a sector for
$N=6\dots18$ and then breaks upward at $N=19$, so it is open too rather than a
clean negative. \rt{obc0} only.
\end{abstract}

\section{The two exponents}

Every rule has two numbers here:
\[
n_\wcc(N)\ \ \text{---\ how many sectors,}\qquad
m(N)=\max_{\text{sector}}\#\{\text{terminal SCCs in it}\}\ \ \text{---\ how many
attractors the busiest sector holds.}
\]
$m\equiv1$ for $243$ of the $256$ rules, so only the thirteen of R9~\S7.3 can
contribute. Nested fragmentation means both grow exponentially.

\begin{figure}[htbp]
\centering
\ifgraphic{../../figures/fig_nested_plane_obc0.pdf}{0.72\textwidth}
\caption{\textbf{The two fragmentation exponents.} Horizontal: growth base of the
sector count. Vertical: growth base of the attractors inside one sector. The
shaded strips are ``exponential''. Eight rules sit in the upper-left band ---
all the fragmentation is inside the sectors --- and two ($29$, $71$) in the
lower-right --- all of it in the sector count. The upper-right corner is nested
fragmentation. \emph{Drawn from $N\le16$ only}: at that size $29$ and $71$ look
linear inside a sector and sit on the lower-right axis, but \S\ref{sec:open}
shows both break upward by $N=19$--$20$, so their vertical positions here are
lower bounds and the corner is not reliably empty.}
\end{figure}

\section{Settled: exponential inside a constant or polynomial sector count}

\begin{center}\small
\begin{tabular}{rlllr}
\toprule
rule & tuple & sector count $n_\wcc$ & attractors in one sector & total base \\
\midrule
$203$, $217$ & \rt{VEII}, \rt{VIEI} & constant, $=1$ at every $N$ & exponential $\psi$ & $\psi$ \\
$219$ & \rt{VEEI} & constant, $=1$ & exponential $\psi$ & $\psi$ \\
$36$ & \rt{IDDV} & constant, $\to10$ & exponential $\psi$ & $\psi$ \\
$44$, $100$ & \rt{IIDV}, \rt{IDIV} & \textbf{polynomial} (linear) & exponential $\psi$ & $\psi$ \\
$104$, $233$ & \rt{DIIV}, \rt{VIIE} & \textbf{polynomial} (linear) & exponential & $1.380278$ \\
\midrule
$29$, $71$ & \rt{EIVD}, \rt{EVID} & \textbf{exponential} & linear then breaks --- \S\ref{sec:open} & $\rho$ \\
$235$, $249$, $37$ & \rt{VEIE}, \rt{VIEE}, \rt{EDDV} & constant, $=1$ & constant ($\le2$) & constant \\
\bottomrule
\end{tabular}
\end{center}

\noindent Rules $203$ and $217$ are the extreme: \emph{one} sector at every $N$,
holding $277$ attractors at $N=16$ and growing like $\psi^N$. The sector
observable returns the constant $1$ however large the system gets, while the
attractor count runs away exponentially --- all of the structure is invisible to
it. Rules $44$ and $100$ are the polynomial version: $n_\wcc=N+3$ exactly for
$100$, with $\psi^N$ attractors inside the largest of those sectors.

\subsection{Three exact constants, all Narayana-like}

The \emph{total} attractor counts obey exact integer recurrences --- found by
recurrence search, not fitted:

\begin{center}\small
\begin{tabular}{llll}
\toprule
rules & recurrence & characteristic root & name \\
\midrule
$203$, $217$, $219$, $36$, $44$, $100$ & $a_N=a_{N-1}+a_{N-3}$ & $x^3=x^2+1$ & supergolden $\psi=1.465571$ \\
$104$, $233$ & $a_N=a_{N-1}+a_{N-4}$ & $x^4=x^3+1$ & $1.380278$ \\
$29$, $71$ & $a_N=a_{N-2}+a_{N-3}$ & $x^3=x+1$ & plastic $\rho=1.324718$ \\
\bottomrule
\end{tabular}
\end{center}

\noindent The first is Narayana's cows sequence
$1,1,1,2,3,4,6,9,13,19,28,41,60,88,129,189,277,406,595$; rules $203/217$ enter it
at $6$, rule $219$ at $3$, and $36/44/100$ at $13$. So six rules share one
sequence at four different offsets.

\subsection{The prettiest fact}

\begin{figure}[htbp]
\centering
\ifgraphic{../../figures/fig_nested_staircase_N16_obc0.pdf}{0.86\textwidth}
\caption{\textbf{How the attractors distribute across sectors} ($N=16$). For
rules $44$ and $100$ the sorted per-sector attractor counts are, exactly, a run
of \emph{consecutive Narayana numbers} --- the same sequence that gives the
total. Rule $36$ has the same total, $595$, but piles $586$ of them into one
sector and leaves nine singletons. Rule $203$ is a single point: one sector,
$277$ attractors. Rule $29$ spreads $86$ attractors thinly over $22$ sectors,
none holding more than $9$.}
\end{figure}

\noindent At $N=16$ rules $44$ and $100$ give
\[
189,\;129,\;88,\;60,\;41,\;28,\;19,\;13,\;9,\;6,\;\dots
\]
which is the Narayana sequence read backwards from its largest term below the
total. Three rules --- $36$, $44$, $100$ --- have the \emph{identical} total of
$595$ attractors at $N=16$ and distribute them completely differently: one giant
sector, or a Narayana staircase.

\section{Open: the fully nested case}
\label{sec:open}

Only $29$ and $71$ have an exponential sector count, so only they can be nested.

\noindent\textbf{Both are undecided, and in the same way.} Each looks exactly
linear inside a sector for as far as one first computes, and each then breaks
upward at the largest $N$ reached. Rule $71$'s in-sector count is
\[
1,2,2,3,3,4,4,5,5,6,6,7,7,\mathbf{9} \qquad (N=6\dots19),
\]
which is $\lceil N/2\rceil-2$ \emph{exactly} for $N=6\dots18$ and then returns
$9$ where the formula says $8$. Its even branch $1,2,3,4,5,6,7$ is still a
perfect line; the break is on the odd branch, at the last point available. Rule
$29$ breaks earlier and further. So the tidy statement ``$71$ is linear, $29$ is
not'' --- which an earlier draft of this report made after stopping at $N=16$ ---
is an artefact of stopping too soon.

\noindent\textbf{Rule 29, quantified.} Its in-sector count oscillates with
parity; the even branch is
\[
3,\;4,\;5,\;6,\;7,\;9,\;12,\;16 \qquad (N=6,8,\dots,20),
\]
which is clearly superlinear --- a line would give $8$ at $N=18$ and $9$ at
$N=20$. But superlinear is not exponential, and the second differences are
$0,0,0,1,1,1$, which is what a \emph{quadratic} looks like. Fitting both:

\begin{center}\small
\begin{tabular}{lrrrr}
\toprule
model & rss & BIC & predicts $N=22$ & predicts $N=26$ \\
\midrule
quadratic in $N$ & $1.571$ & $-6.78$ & $19.2$ & $28.1$ \\
exponential, base $1.1206$ & $1.672$ & $-8.37$ & $18.8$ & $29.7$ \\
\bottomrule
\end{tabular}
\end{center}

\begin{figure}[htbp]
\centering
\ifgraphic{../../figures/fig_nested_rule29_obc0.pdf}{0.72\textwidth}
\caption{\textbf{The undecided case.} Rule $29$'s attractors in the largest
sector, even $N$, with both models. They agree to within $6\%$ everywhere data
exists and separate by less than one attractor at $N=22$.}
\end{figure}

\noindent A BIC gap of $1.6$ is nothing, and the two models differ by $0.4$
attractors at $N=22$ and $1.6$ at $N=26$. Settling it needs $N\gtrsim26$, i.e.
$6.7\times10^{7}$ states --- feasible only with a bitwise successor and a
compact union-find, not the present path. \textbf{Verdict: open.} The
automated classifier calls it exponential (base $1.111$) and that verdict is
recorded, but it rests on a short, parity-oscillating series and should not be
quoted as a result.

\section{What this says}

The eight settled rules answer the second half of the question affirmatively and
sharply: \emph{a constant or polynomial sector count is fully compatible with
exponentially many attractors inside a single sector.} Rules $203$ and $217$ are
the cleanest statement of it --- one sector, forever, containing $\psi^N$
attractors --- and they are the reason R9~\S7.3 insists that counting sectors is
not counting attractors.

Whether a rule can do \emph{both} at once remains open on \emph{two} rules, $29$
and $71$, by a margin of about one attractor at the largest size computed here.
Both are linear inside a sector for a long stretch and both break upward at the
end of it, which is the least conclusive shape a finite series can have: it is
equally the onset of a slow exponential and a finite-size wobble on a line.

\section*{Reproduce}

\begin{verbatim}
python -m qca_fragmentation.viz.nested --bc obc0 --rebuild
\end{verbatim}

\noindent writes \fn{analytics/nested\_fragmentation\_obc0.json} and the three
figures. The module is \fn{code/qca\_fragmentation/viz/nested.py};
\fn{rule29\_model\_comparison()} reproduces the table of \S\ref{sec:open}.

\end{document}
